\begin{table}[!htbp]
\centering
\caption{Local support and fuzzy-heterogeneity identification diagnostics}
\label{tab:hte-2017_ccpp-diagnostics}
\small
\resizebox{\textwidth}{!}{%
\begin{tabular}{lrrrrrll}
\toprule
Moderator & CCPP left & CCPP right & Min. cell & Min. SW \(F\) & KP \(F\) & Support & IV gate \\
\midrule
Log baseline CCPP population & 38 & 21 & 21 & 11.78 & 2.04 & Pass & Pass \\
District-capital CCPP & 38 & 21 & 2 & -- & -- & Fail & Fail \\
Baseline CCPP deprivation & 37 & 21 & 21 & 28.54 & 12.12 & Pass & Pass \\
Distance to corresponding district capital & 38 & 21 & 21 & 55.66 & 4.24 & Pass & Pass \\
Baseline CCPP economic development & 38 & 21 & 21 & 5.01 & 2.06 & Pass & Fail \\
CCPP altitude & 38 & 21 & 21 & 5.92 & 3.15 & Pass & Fail \\
\bottomrule
\end{tabular}}
\parbox{0.97\linewidth}{\footnotesize \textit{Notes:} Diagnostics come from the pooled, fully interacted local-linear 2SLS model in the fixed \(h=0.0075\) window. Each RUV community contributes one observation before triangular kernel weighting; primary inference clusters by district. CCPP left and right are unique assignment clusters. For binary moderators, Min. cell is the smallest moderator-by-side CCPP count; for continuous moderators it is the smaller side count. The causal gate requires support, underidentification rejection at five percent, and the minimum Sanderson--Windmeijer conditional \(F\) strictly above 10. The Kleibergen--Paap statistic is supplementary. Source: RUV, CMAN, and INEI-assisted Census 2017.}
\end{table}
